
\documentclass{article}
\usepackage{colortbl}
\usepackage{makecell}
\usepackage{multirow}
\usepackage{supertabular}

\begin{document}

\newcounter{utterance}

\twocolumn

{ \footnotesize  \setcounter{utterance}{1}
\setlength{\tabcolsep}{0pt}
\begin{supertabular}{c@{$\;$}|p{.15\linewidth}@{}p{.15\linewidth}p{.15\linewidth}p{.15\linewidth}p{.15\linewidth}p{.15\linewidth}}

    \# & $\;$A & \multicolumn{4}{c}{Game Master} & $\;\:$B\\
    \hline 

    \theutterance \stepcounter{utterance}  

    & & \multicolumn{4}{p{0.6\linewidth}}{\cellcolor[rgb]{0.9,0.9,0.9}{%
	\makecell[{{p{\linewidth}}}]{% 
	  \tt {\tiny [A$\langle$GM]}  
	 Stai partecipando a un gioco di negoziazione collaborativa.\\ \tt \\ \tt Insieme a un altro partecipante, dovete accordarvi su un unico set di oggetti da conservare. Ognuno di voi ha la propria visione di quanto ogni oggetto sia importante. Non sai come l'altro partecipante valuta gli oggetti. Inoltre, ti viene fornito lo sforzo che ogni oggetto richiede.\\ \tt Potete accordarvi su un set solo se lo sforzo totale degli oggetti selezionati non supera un limite condiviso:\\ \tt \\ \tt LIMITE = 2984\\ \tt \\ \tt Ecco i valori di sforzo dei singoli oggetti:\\ \tt \\ \tt Sforzo oggetto = {"A21": 350, "B93": 276, "A11": 616, "B33": 737, "C08": 531, "B39": 389, "B96": 24, "C85": 125, "A82": 338, "B31": 356, "B20": 143, "A60": 117, "B09": 257, "C81": 921, "B21": 789}\\ \tt \\ \tt Ecco la tua visione personale sull'importanza di ogni oggetto:\\ \tt \\ \tt Valori di importanza dell'oggetto = {"A21": 138, "B93": 583, "A11": 868, "B33": 822, "C08": 783, "B39": 65, "B96": 262, "C85": 121, "A82": 508, "B31": 780, "B20": 461, "A60": 484, "B09": 668, "C81": 389, "B21": 808}\\ \tt \\ \tt Obiettivo:\\ \tt \\ \tt Il tuo obiettivo è negoziare un set condiviso di oggetti che ti avvantaggi il più possibile (cioè, massimizzi l'importanza totale per TE), rimanendo entro il LIMITEE. Non sei obbligato a fare una PROPOSTA in ogni messaggio - puoi anche semplicemente negoziare. Tutte le tattiche sono permesse!\\ \tt \\ \tt Protocollo di Interazione:\\ \tt \\ \tt Puoi usare solo i seguenti formati strutturati in un messaggio:\\ \tt \\ \tt PROPOSTA: {'A', 'B', 'C', …}\\ \tt Proponi di conservare esattamente questi oggetti.\\ \tt RIFIUTO: {'A', 'B', 'C', …}\\ \tt Rifiuta esplicitamente la proposta dell'avversario.\\ \tt ARGOMENTAZIONE: {'...'}\\ \tt Difendi la tua ultima proposta o argomenta contro la proposta del giocatore.\\ \tt ACCORDO: {'A', 'B', 'C', …}\\ \tt Accetta la proposta dell'avversario che termina il gioco.\\ \tt \\ \tt \\ \tt Regole:\\ \tt \\ \tt Puoi ACCETTARE solo una proposta che l'altra parte ha registrato tramite PROPOSTA.\\ \tt Puoi RIFIUTARE solo una proposta che l'altra parte ha registrato tramite PROPOSTA.\\ \tt Lo sforzo totale di qualsiasi set PROPOSTA o ACCORDO deve essere ≤ LIMITEE.\\ \tt NON rivelare i tuoi punteggi di importanza nascosti.\\ \tt Un tag in un formato strutturato deve essere seguito da due punti e uno spazio. L'argomento deve essere un set python contenente 0 o più stringhe.\\ \tt Quindi, deve essere della forma TAG: {...}\\ \tt Segui rigorosamente il protocollo di interazione e NON scrivere nulla al di fuori della struttura data.\\ \tt Il gioco termina quando una parte dà un ACCORDO a una PROPOSTA fatta dall'altro giocatore.\\ \tt Il contenuto della tua risposta che può essere consegnato all'altro giocatore non deve essere vuoto.\\ \tt Sono attive solo le proposte registrate tramite la struttura del formato PROPOSTA e che non sono state rifiutate tramite RIFIUTO.\\ \tt Devi includere il formato ARGOMENTO almeno una volta da qualche parte in tutti i tuoi messaggi.\\ \tt  
	  } 
	   } 
	   } 
	 & \\ 
 

    \theutterance \stepcounter{utterance}  

    & \multicolumn{4}{p{0.6\linewidth}}{\cellcolor[rgb]{0.8,1,0.9}{%
	\makecell[{{p{\linewidth}}}]{% 
	  \tt {\tiny [A$\rangle$GM]}  
	 PROPOSTA: {'A11', 'B33', 'C08', 'B31', 'B21'} \\ \tt ARGOMENTAZIONE: {'Ho selezionato questi oggetti perché ritengo che siano fondamentali per raggiungere un buon equilibrio tra sforzo e importanza, garantendo al contempo che il limite di sforzo totale non venga superato. Sono disposto a discutere e negoziare per trovare un accordo che soddisfi entrambe le parti.'} 
	  } 
	   } 
	   } 
	 & & \\ 
 

    \theutterance \stepcounter{utterance}  

    & & \multicolumn{4}{p{0.6\linewidth}}{\cellcolor[rgb]{0.9,0.9,0.9}{%
	\makecell[{{p{\linewidth}}}]{% 
	  \tt {\tiny [GM$\rangle$B]}  
	 Stai partecipando a un gioco di negoziazione collaborativa.\\ \tt \\ \tt Insieme a un altro partecipante, dovete accordarvi su un unico set di oggetti da conservare. Ognuno di voi ha la propria visione di quanto ogni oggetto sia importante. Non sai come l'altro partecipante valuta gli oggetti. Inoltre, ti viene fornito lo sforzo che ogni oggetto richiede.\\ \tt Potete accordarvi su un set solo se lo sforzo totale degli oggetti selezionati non supera un limite condiviso:\\ \tt \\ \tt LIMITE = 2984\\ \tt \\ \tt Ecco i valori di sforzo dei singoli oggetti:\\ \tt \\ \tt Sforzo oggetto = {"A21": 350, "B93": 276, "A11": 616, "B33": 737, "C08": 531, "B39": 389, "B96": 24, "C85": 125, "A82": 338, "B31": 356, "B20": 143, "A60": 117, "B09": 257, "C81": 921, "B21": 789}\\ \tt \\ \tt Ecco la tua visione personale sull'importanza di ogni oggetto:\\ \tt \\ \tt Valori di importanza dell'oggetto = {"A11": 65, "B33": 121, "B21": 138, "C08": 262, "B31": 389, "B09": 461, "B93": 484, "A82": 508, "A60": 583, "B20": 668, "C81": 780, "B96": 783, "A21": 808, "C85": 822, "B39": 868}\\ \tt \\ \tt Obiettivo:\\ \tt \\ \tt Il tuo obiettivo è negoziare un set condiviso di oggetti che ti avvantaggi il più possibile (cioè, massimizzi l'importanza totale per TE), rimanendo entro il LIMITEE. Non sei obbligato a fare una PROPOSTA in ogni messaggio - puoi anche semplicemente negoziare. Tutte le tattiche sono permesse!\\ \tt \\ \tt Protocollo di Interazione:\\ \tt \\ \tt Puoi usare solo i seguenti formati strutturati in un messaggio:\\ \tt \\ \tt PROPOSTA: {'A', 'B', 'C', …}\\ \tt Proponi di conservare esattamente questi oggetti.\\ \tt RIFIUTO: {'A', 'B', 'C', …}\\ \tt Rifiuta esplicitamente la proposta dell'avversario.\\ \tt ARGOMENTAZIONE: {'...'}\\ \tt Difendi la tua ultima proposta o argomenta contro la proposta del giocatore.\\ \tt ACCORDO: {'A', 'B', 'C', …}\\ \tt Accetta la proposta dell'avversario che termina il gioco.\\ \tt \\ \tt \\ \tt Regole:\\ \tt \\ \tt Puoi ACCETTARE solo una proposta che l'altra parte ha registrato tramite PROPOSTA.\\ \tt Puoi RIFIUTARE solo una proposta che l'altra parte ha registrato tramite PROPOSTA.\\ \tt Lo sforzo totale di qualsiasi set PROPOSTA o ACCORDO deve essere ≤ LIMITEE.\\ \tt NON rivelare i tuoi punteggi di importanza nascosti.\\ \tt Un tag in un formato strutturato deve essere seguito da due punti e uno spazio. L'argomento deve essere un set python contenente 0 o più stringhe.\\ \tt Quindi, deve essere della forma TAG: {...}\\ \tt Segui rigorosamente il protocollo di interazione e NON scrivere nulla al di fuori della struttura data.\\ \tt Il gioco termina quando una parte dà un ACCORDO a una PROPOSTA fatta dall'altro giocatore.\\ \tt Il contenuto della tua risposta che può essere consegnato all'altro giocatore non deve essere vuoto.\\ \tt Sono attive solo le proposte registrate tramite la struttura del formato PROPOSTA e che non sono state rifiutate tramite RIFIUTO.\\ \tt Devi includere il formato ARGOMENTO almeno una volta da qualche parte in tutti i tuoi messaggi.\\ \tt \\ \tt \\ \tt Ora riceverai il primo messaggio dell'altro giocatore.\\ \tt \\ \tt PROPOSTA: {'A11', 'B33', 'C08', 'B31', 'B21'} \\ \tt ARGOMENTAZIONE: {'Ho selezionato questi oggetti perché ritengo che siano fondamentali per raggiungere un buon equilibrio tra sforzo e importanza, garantendo al contempo che il limite di sforzo totale non venga superato. Sono disposto a discutere e negoziare per trovare un accordo che soddisfi entrambe le parti.'} 
	  } 
	   } 
	   } 
	 & \\ 
 

    \theutterance \stepcounter{utterance}  

    & & & \multicolumn{4}{p{0.6\linewidth}}{\cellcolor[rgb]{1,0.85,0.72}{%
	\makecell[{{p{\linewidth}}}]{% 
	  \tt {\tiny [GM$\langle$B]}  
	 RIFIUTO: {'A11', 'B33', 'C08', 'B31', 'B21'}\\ \tt ARGOMENTAZIONE: {'La proposta iniziale sembra non tenere conto di alcuni oggetti di importanza critica. Sono preoccupato che escludendo oggetti come "A21", "B93" e "C85", potremmo perdere oportunità significative di aumentare l'importanza totale senza superare il limite di sforzo. Sarebbe utile discutere ulteriormente le priorità e valutare se ci sono compromessi possibili che potrebbero inclusi altri oggetti di alto valore.'} 
	  } 
	   } 
	   } 
	  \\ 
 

    \theutterance \stepcounter{utterance}  

    & & \multicolumn{4}{p{0.6\linewidth}}{\cellcolor[rgb]{0.9,0.9,0.9}{%
	\makecell[{{p{\linewidth}}}]{% 
	  \tt {\tiny [A$\langle$GM]}  
	 RIFIUTO: {'A11', 'B33', 'C08', 'B31', 'B21'}\\ \tt ARGOMENTAZIONE: {'La proposta iniziale sembra non tenere conto di alcuni oggetti di importanza critica. Sono preoccupato che escludendo oggetti come "A21", "B93" e "C85", potremmo perdere oportunità significative di aumentare l'importanza totale senza superare il limite di sforzo. Sarebbe utile discutere ulteriormente le priorità e valutare se ci sono compromessi possibili che potrebbero inclusi altri oggetti di alto valore.'} 
	  } 
	   } 
	   } 
	 & \\ 
 

    \theutterance \stepcounter{utterance}  

    & \multicolumn{4}{p{0.6\linewidth}}{\cellcolor[rgb]{0.8,1,0.9}{%
	\makecell[{{p{\linewidth}}}]{% 
	  \tt {\tiny [A$\rangle$GM]}  
	 ARGOMENTAZIONE: {'Capisco le preoccupazioni riguardo alla mia proposta iniziale e sono aperto a esplorare alternative. Tuttavia, vorrei sottolineare che gli oggetti come "A11", "B33", "C08", "B31" e "B21" rappresentano un buon punto di partenza a causa del loro equilibrio tra sforzo e importanza. Forse potremmo considerare di sostituire alcuni di questi oggetti con quelli che ritieni più critici, come "A21", "B93" e "C85", pur mantenendo il limite di sforzo totale. Sarebbe utile conoscere la tua visione su come questi oggetti potrebbero essere integrati in una proposta che massimizzi il valore totale.'} \\ \tt PROPOSTA: {'A21', 'B93', 'B31', 'B21', 'A11'} 
	  } 
	   } 
	   } 
	 & & \\ 
 

    \theutterance \stepcounter{utterance}  

    & & \multicolumn{4}{p{0.6\linewidth}}{\cellcolor[rgb]{0.9,0.9,0.9}{%
	\makecell[{{p{\linewidth}}}]{% 
	  \tt {\tiny [GM$\rangle$B]}  
	 ARGOMENTAZIONE: {'Capisco le preoccupazioni riguardo alla mia proposta iniziale e sono aperto a esplorare alternative. Tuttavia, vorrei sottolineare che gli oggetti come "A11", "B33", "C08", "B31" e "B21" rappresentano un buon punto di partenza a causa del loro equilibrio tra sforzo e importanza. Forse potremmo considerare di sostituire alcuni di questi oggetti con quelli che ritieni più critici, come "A21", "B93" e "C85", pur mantenendo il limite di sforzo totale. Sarebbe utile conoscere la tua visione su come questi oggetti potrebbero essere integrati in una proposta che massimizzi il valore totale.'} \\ \tt PROPOSTA: {'A21', 'B93', 'B31', 'B21', 'A11'} 
	  } 
	   } 
	   } 
	 & \\ 
 

    \theutterance \stepcounter{utterance}  

    & & & \multicolumn{4}{p{0.6\linewidth}}{\cellcolor[rgb]{1,0.85,0.72}{%
	\makecell[{{p{\linewidth}}}]{% 
	  \tt {\tiny [GM$\langle$B]}  
	 ARGOMENTAZIONE: {'Apprezzo la tua disponibilità a considerare alternative e a integrare alcuni degli oggetti che ritengo più critici. La tua nuova proposta include "A21" e "B93", che sono tra quelli che considero di alta importanza. Tuttavia, noto che sono stati esclusi "C08" e "C85", che potrebbero avere un impatto significativo sul valore totale. Sarebbe interessante valutare come l'inclusione di "C85" potrebbe influenzare la proposta, considerando il suo potenziale valore. Inoltre, proporrei di valutare se "A60" e "B20" potrebbero essere più efficaci di "A11" in termini di equilibrio tra sforzo e importanza, mantenendo comunque il limite di sforzo totale.'}\\ \tt RIFIUTO: {'A21', 'B93', 'B31', 'B21', 'A11'}\\ \tt PROPOSTA: {'A21', 'B93', 'C85', 'B31', 'B96'} 
	  } 
	   } 
	   } 
	  \\ 
 

    \theutterance \stepcounter{utterance}  

    & & \multicolumn{4}{p{0.6\linewidth}}{\cellcolor[rgb]{0.9,0.9,0.9}{%
	\makecell[{{p{\linewidth}}}]{% 
	  \tt {\tiny [A$\langle$GM]}  
	 ARGOMENTAZIONE: {'Apprezzo la tua disponibilità a considerare alternative e a integrare alcuni degli oggetti che ritengo più critici. La tua nuova proposta include "A21" e "B93", che sono tra quelli che considero di alta importanza. Tuttavia, noto che sono stati esclusi "C08" e "C85", che potrebbero avere un impatto significativo sul valore totale. Sarebbe interessante valutare come l'inclusione di "C85" potrebbe influenzare la proposta, considerando il suo potenziale valore. Inoltre, proporrei di valutare se "A60" e "B20" potrebbero essere più efficaci di "A11" in termini di equilibrio tra sforzo e importanza, mantenendo comunque il limite di sforzo totale.'}\\ \tt RIFIUTO: {'A21', 'B93', 'B31', 'B21', 'A11'}\\ \tt PROPOSTA: {'A21', 'B93', 'C85', 'B31', 'B96'} 
	  } 
	   } 
	   } 
	 & \\ 
 

    \theutterance \stepcounter{utterance}  

    & \multicolumn{4}{p{0.6\linewidth}}{\cellcolor[rgb]{0.8,1,0.9}{%
	\makecell[{{p{\linewidth}}}]{% 
	  \tt {\tiny [A$\rangle$GM]}  
	 ARGOMENTAZIONE: {'Apprezzo la tua proposta e la tua analisi dettagliata. L'inclusione di "C85" e "B96" è interessante, poiché porterà una buona quantità di importanza con un sforzo relativamente contenuto. Tuttavia, sono preoccupato che l'esclusione di "B21" potrebbe diminuire significativamente il valore totale della proposta. "B21" ha un alto valore per me, e la sua rimozione potrebbe influire negativamente sul nostro obiettivo di massimizzare l'importanza totale. Vorrei proporre un compromesso: manteniamo "B21" e "B31" per la loro alta importanza, e includiamo "C85" per il suo valore aggiunto, ma sostituiamo "B96" con "A60" per ottenere un miglior equilibrio tra sforzo e importanza.'} \\ \tt RIFIUTO: {'A21', 'B93', 'C85', 'B31', 'B96'}\\ \tt PROPOSTA: {'A21', 'B93', 'C85', 'B21', 'B31'} 
	  } 
	   } 
	   } 
	 & & \\ 
 

    \theutterance \stepcounter{utterance}  

    & & \multicolumn{4}{p{0.6\linewidth}}{\cellcolor[rgb]{0.9,0.9,0.9}{%
	\makecell[{{p{\linewidth}}}]{% 
	  \tt {\tiny [GM$\rangle$B]}  
	 ARGOMENTAZIONE: {'Apprezzo la tua proposta e la tua analisi dettagliata. L'inclusione di "C85" e "B96" è interessante, poiché porterà una buona quantità di importanza con un sforzo relativamente contenuto. Tuttavia, sono preoccupato che l'esclusione di "B21" potrebbe diminuire significativamente il valore totale della proposta. "B21" ha un alto valore per me, e la sua rimozione potrebbe influire negativamente sul nostro obiettivo di massimizzare l'importanza totale. Vorrei proporre un compromesso: manteniamo "B21" e "B31" per la loro alta importanza, e includiamo "C85" per il suo valore aggiunto, ma sostituiamo "B96" con "A60" per ottenere un miglior equilibrio tra sforzo e importanza.'} \\ \tt RIFIUTO: {'A21', 'B93', 'C85', 'B31', 'B96'}\\ \tt PROPOSTA: {'A21', 'B93', 'C85', 'B21', 'B31'} 
	  } 
	   } 
	   } 
	 & \\ 
 

    \theutterance \stepcounter{utterance}  

    & & & \multicolumn{4}{p{0.6\linewidth}}{\cellcolor[rgb]{1,0.85,0.72}{%
	\makecell[{{p{\linewidth}}}]{% 
	  \tt {\tiny [GM$\langle$B]}  
	 ARGOMENTAZIONE: {'La tua controproposta è interessante e mostra una buona comprensione delle priorità. Includere "B21" e "B31" mantiene un livello di importanza elevato, mentre "C85" aggiunge ulteriore valore. Sono d'accordo con la tua valutazione che "A60" potrebbe essere una buona scelta per migliorare l'equilibrio, ma nella tua proposta finale, hai sostituito "A60" con "A21" e mantenuto "B93". Sono curioso di capire il motivo dietro questa scelta, poiché "B93" ha un sforzo relativamente alto e potrebbe potenzialmente essere sostituito con un oggetto che offre un miglior rapporto sforzo-importanza. Tuttavia, la tua proposta attuale sembra essere una buona base per un accordo, considerando che massimizza l'importanza totale senza superare il limite di sforzo.'}\\ \tt ACCORDO: {'A21', 'B93', 'C85', 'B21', 'B31'} 
	  } 
	   } 
	   } 
	  \\ 
 

    \theutterance \stepcounter{utterance}  

    & & & \multicolumn{2}{p{0.3\linewidth}}{\cellcolor[rgb]{0.95,0.95,0.95}{%
	\makecell[{{p{\linewidth}}}]{% 
	  \tt {\tiny [GM$|$GM]}  
	 game successful 
	  } 
	   } 
	   } 
	 & & \\ 
 

    \theutterance \stepcounter{utterance}  

    & & & \multicolumn{2}{p{0.3\linewidth}}{\cellcolor[rgb]{0.95,0.95,0.95}{%
	\makecell[{{p{\linewidth}}}]{% 
	  \tt {\tiny [GM$|$GM]}  
	 end game 
	  } 
	   } 
	   } 
	 & & \\ 
 

\end{supertabular}
}

\end{document}
